\documentclass[final,hyperref={pdfpagelabels=false}]{beamer}
\mode<presentation>{\usetheme{I6pd2}}
\usepackage[english]{babel}
\usepackage{amsmath,amsthm, amssymb, latexsym}
\usepackage[orientation=portrait,size=a0,debug]{beamerposter}
\usepackage{xcolor,colortbl}
\usepackage{tipa}
\usepackage{natbib}
\usepackage{multirow}
\usepackage{subfig}
\usepackage{tikz}
\usetikzlibrary{calc}
\usetikzlibrary{scopes}

%\usepackage{grffile}
%\usefonttheme[onlymath]{serif}
%\boldmath
% change list indention level
% \setdefaultleftmargin{3em}{}{}{}{}{}

\bibpunct[:]{(}{)}{,}{a}{}{,}
\newcommand{\ipa}[1]{\textipa{#1}} 

\usepackage{snapshot} % will write a .dep file with all dependencies, allows for easy bundling

\usepackage{array,booktabs,tabularx}
\newcolumntype{Z}{>{\centering\arraybackslash}X} % centered tabularx columns
\newcommand{\pphantom}{\textcolor{ta3aluminium}} % phantom introduces a vertical space in p formatted table columns??!!
\listfiles

%%%%
%%%%%%
\makeatletter
\define@key{beamerframe}{t}[true]{% top
  \beamer@frametopskip=.2cm plus .5\paperheight\relax%
  \beamer@framebottomskip=0pt plus 1fill\relax%
  \beamer@frametopskipautobreak=\beamer@frametopskip\relax%
  \beamer@framebottomskipautobreak=\beamer@framebottomskip\relax%
% \def\beamer@initfirstlineunskip{%
%   \def\beamer@firstlineitemizeunskip{%
%     \vskip-\partopsep\vskip-\topsep\vskip-\parskip%
%     \global\let\beamer@firstlineitemizeunskip=\relax}%
%   \everypar{\global\let\beamer@firstlineitemizeunskip=\relax}}
  \def\beamer@initfirstlineunskip{}%
}
\makeatother

\graphicspath{{figures/}}
 
\title{\huge Title}
\author{Pertti Palo$^{\mathrm{a}}$ and Steven Lulich$^{\mathrm{b}}$}
\newcommand{\emails}{$^{\mathrm{a}}$pertti.palo@taurlin.org, $^{\mathrm{b}}$slulich@indiana.edu}
\institute{Department of Speech, Language and Hearing Sciences, Indiana University, Bloomington, IN}
\date{6th October 2021}

\newlength{\columnheight}
\setlength{\columnheight}{105cm}

\newcommand{\kommentti}[1]{{\bf[#1]}}


\begin{document}

\begin{frame}{} 
  \vspace{-1cm}
  \begin{columns}[t]
    \begin{column}{.45\linewidth}
      
      \begin{block}{Introduction: What do we know about acoustic
          reaction times?}

        \begin{columns}
          \hspace*{-.25cm}
          \begin{column}{.3\textwidth}
            \begin{itemize}
            \item \cite{RastleEtAl-CharacterizingMotorExecution-2005} measured
            delayed naming reaction times.
            \item They report Acoustic RTs and onset durations (OD) for all
              phonotactically legal English consonant onsets.
            \item The inverse correlation pattern was first discovered by
              \cite{Fowler-PerceptualCentersSpeech-1979}.
            \end{itemize}
          \end{column}
          \begin{column}{.6\textwidth}
            \kommentti{Rastle figure here.}
            % \vspace*{.5cm}
            % \input{figures/Rastle_ICPhS.tex}
          \end{column}
        \end{columns}
      \end{block}

      \begin{block}{Materials: The Experiment} 

        'Simple' task of finding acoustic and movement onsets in
        simultaneously recorded sound data and articulatory data from
        Ultrasound Tongue Imaging (UTI).

        /CVC/ and /VC/ English lexical words produced by one
        male speaker: {\em at, eat, ought, back, beat, bought, dat, deep,
          dot, fat, feet, fought, gap, geek, got, hat, heat, hot, cat,
          keep, caught, lack, leap, lot, map, meet, mock, nat, neat, not,
          pack, Pete, pop, rat, reap, rock, sat, seat, sought, shack,
          sheet, shop, tap, teak, talk, whack, wheat,} and {\em what}.
      \end{block}
      
      \begin{block}{Methods: Recording and annotation}\footnotesize
        \begin{itemize}
        \item \kommentti{Say something about methods.}
        \end{itemize}
      \end{block}

      \begin{block}{Results: Acoustic RTs}
        \begin{itemize}
        \item Acoustic RT (AcRT) behaves as expected
          \citep{RastleEtAl-CharacterizingMotorExecution-2005,
          MooshammerEtAl-BridgingPlanningExecution-2012}.
        \item Statistical analysis yields an expected result
          \citep{Fowler-PerceptualCentersSpeech-1979,
          RastleEtAl-CharacterizingMotorExecution-2005}:
          AcRT $\sim -\frac{1}{2}$ OD.
        \end{itemize}
      \end{block}

    \end{column}
    \begin{column}{.45\linewidth}
          
      \begin{block}{Results: Articulatory to Acoustic Interval (AAI)}
        \begin{itemize}
        \item Again statistical analysis yields: AAI $\sim
          -\frac{1}{2}$ OD.
        \item However, a better predictor for AAI than OD is phoneme.
        \end{itemize}
        % \vspace*{.5cm}
        % \begin{center}
        %   \input{figures/OD_vs_AAI.tex}
        % \end{center}
      \end{block}

      \begin{block}{Discussion: Future work}

        \begin{itemize}
        \item There is no data like more data.
        \end{itemize}
      \end{block}

    \end{column}
  \end{columns}
      
  \vspace*{1cm}
  
  \begin{columns}
    \begin{column}{.945\linewidth}

      \begin{block}{References}

        \footnotesize
        % This needs to be replaced by a symlink or a fresh copy.
        \bibliography{main.bib}
        \bibliographystyle{apalike}

      \end{block}

      \vspace*{1cm}

      \begin{block}{Acknowledgements}
        %\footnotesize
        \begin{minipage}{0.5\textwidth}
          % \includegraphics[scale=1]{figures/rbo.jpg}
          We wish to thank Steve Cowen for assistance with the
          ultrasound recordings and Professor Alan Wrench for advice
          and help with using AAA.
        \end{minipage}
        % \begin{minipage}{0.5\textwidth}
        %   \hspace{10cm}\includegraphics[height=3.5cm]{figures/CASL_Sub_logo.jpg} 
        % \end{minipage}
      \end{block}
    \end{column}
  \end{columns}

  
\end{frame}


\end{document}

        % {Connection with speech cycling}
        % \begin{center}
        %   \includegraphics[width=.7\columnwidth]{Fowler_cycling}
        % \end{center}
        % \begin{itemize}
        % \item Figure from \cite{Fowler:NMC:1981}.
        % \item Metronome click at 0 ms.
        % \item Relevant intervals are: 1. silence between words, 3. onset
        %   consonant, and 4. vowel.
        % \end{itemize}

        % {Future work and questions}


%%%%%%%%%%%%%%%%%%%%%%%%%%%%%%%%%%%%%%%%%%%%%%%%%%%%%%%%%%%%%%%%%%%%%%%%%%%%%%%%%%%%%%%%%%%%%%%%%%%%
%%% Local Variables: 
%%% mode: latex
%%% TeX-PDF-mode: t
%%% End:
